% ============================================================
%  Supplementary Information
%  Spatial geometry of disaster displacement is linked to
%  population recovery speed
% ============================================================
\documentclass[11pt]{article}

\usepackage[utf8]{inputenc}
\usepackage[T1]{fontenc}
\usepackage{mathptmx}
\usepackage[margin=2.5cm]{geometry}
\usepackage{amsmath,amssymb}
\usepackage{graphicx}
\usepackage{booktabs}
\usepackage{natbib}
\usepackage{hyperref}
\usepackage[font=small,labelfont=bf]{caption}
\usepackage{array}
\usepackage{xcolor}

% ---- macros (same as main.tex) ----
\newcommand{\dnear}{\delta_{\mathrm{near}}}
\newcommand{\dpeak}{D_{\mathrm{peak}}}
\newcommand{\apred}{\alpha_{\mathrm{pred}}}
\newcommand{\aemp}{\alpha_{\mathrm{emp}}}
\newcommand{\Ds}{D_s}
\newcommand{\Dnorm}{D_{\mathrm{norm}}}
\newcommand{\rhosp}{\rho}
\newcommand{\aunit}{\alpha_{\mathrm{unit}}}

\title{Supplementary Information\\[4pt]
Spatial geometry of disaster displacement governs\\
population recovery speed}
\author{}
\date{}

\begin{document}
\maketitle

\setcounter{section}{0}
\renewcommand{\thesection}{S\arabic{section}}
\renewcommand{\thetable}{S\arabic{table}}
\renewcommand{\thefigure}{S\arabic{figure}}

% ============================================================
%  S1  Event Catalogue
% ============================================================
\section{Event catalogue and sample selection}
\label{sec:si_events}

The screened cohort comprises 18~events, and all 18 enter the
primary event-level analysis. Table~\ref{tab:events} lists this
full 18-event event-level sample, ordered by hazard type and then by
descending~$\aemp$. Events were selected from the
Facebook Disaster Maps archive using
the following quality criteria: (i) at least three consecutive
post-peak time steps showing monotone decay ($n_{\mathrm{mono}} \geq 3$),
so that a power-law slope can be estimated; (ii) a clearly identifiable
peak in $D(t)$ with no ambiguity arising from multiple activation
periods or overlapping events; and (iii) sufficient tile coverage within
the 200-km analysis disk ($\geq 3$ populated tiles per annulus at peak).
Duplicate activations of the same physical event and activations with
fewer than two near-field tiles were excluded. Two events (the eastern
Bangladesh floods and the Boise County wildfire) yield no unit-level
fits and therefore do not enter the geo-unit analyses.

\begin{table}[h]
\centering
\caption{Event catalogue.  $\dpeak$: aggregated peak displacement
(dimensionless population ratio deviation).  $\aemp$: empirical
log--log decay rate (monotone-truncation method, window
$[t_{\mathrm{start}},\,t_{\mathrm{end}}]$~h).
$\dnear$: mean near-field displacement at peak time.
$D_\infty$: normalised residual displacement.}
\label{tab:events}
\small
\setlength{\tabcolsep}{4pt}
\begin{tabular}{@{}llllrrrrcr@{}}
\toprule
Event & Country & Hazard & Type
  & $\dpeak$ & $\aemp$ & $R^2$ & $\dnear$
  & Window (h) & $D_\infty$ \\
\midrule
\multicolumn{10}{l}{\textit{Earthquake}} \\[2pt]
2023 T\"{u}rkiye--Syria EQ    & T\"{u}rkiye & EQ & INFL & 0.401 & 0.223 & 0.959 &  0.057 & 24--96  & 0.569 \\
\midrule
\multicolumn{10}{l}{\textit{Tropical cyclone}} \\[2pt]
Hurricane Beryl (Texas, US)       & USA         & TC & EVAC & 0.241 & 0.800 & 0.990 & $-0.279$ & 24--144 & 0.117 \\
Hurricane Beryl (QR, Mexico)      & Mexico      & TC & EVAC & 0.236 & 0.741 & 0.979 & $-0.442$ & 24--144 & 0.083 \\
Tropical Storm Kristine (PH)      & Philippines & TC & EVAC & 0.297 & 0.711 & 0.943 & $-0.041$ & 24--96  & 0.500 \\
Typhoon Yagi (N.~Vietnam)         & Vietnam     & TC & EVAC & 0.317 & 0.522 & 0.955 & $-0.547$ & 24--120 & 0.283 \\
Tropical Storm Yagi (Philippines) & Philippines & TC & EVAC & 0.139 & 0.467 & 0.934 & $-0.227$ & 24--192 & 0.717 \\
Hurricane Beryl (Pre-landfall)    & Jamaica     & TC & EVAC & 0.564 & 0.387 & 0.894 & $-0.404$ & 24--216 & 0.290 \\
Typhoon Krathon (Taiwan)          & Taiwan      & TC & INFL & 0.134 & 0.101 & 0.949 &  0.111 & 24--96  & 0.640 \\
\midrule
\multicolumn{10}{l}{\textit{Flood}} \\[2pt]
2024 Nepal Floods (Bagmati/Koshi) & Nepal      & FL & EVAC & 0.212 & 1.215 & 0.944 & $-0.120$ & 24--120 & 0.434 \\
2024 E.~Bangladesh Floods         & Bangladesh & FL & INFL & 0.097 & 0.472 & 0.940 &  0.028 & 24--120 & 0.454 \\
2024 Rio Grande do Sul Floods     & Brazil     & FL & NEUT & 0.109 & 0.390 & 0.932 &  0.001 & 24--72  & 0.209 \\
2024 Gujarat Floods               & India      & FL & EVAC & 0.139 & 0.159 & 0.832 & $-0.133$ & 24--144 & 0.655 \\
2024 Moldova Floods               & Moldova    & FL & EVAC & 0.054 & 0.071 & 0.393 & $-0.029$ & 24--72  & 0.693 \\
2024 Spain Floods (DANA)          & Spain      & FL & INFL & 0.127 & 0.058 & 0.813 &  0.203 & 24--96  & 0.864 \\
2024 C.--E.~Europe Floods         & Poland     & FL & INFL & 0.143 & $-0.015$ & 0.141 &  0.101 & 24--144 & 0.872 \\
\midrule
\multicolumn{10}{l}{\textit{Wildfire}} \\[2pt]
2024 Quito Wildfires              & Ecuador    & WF & NEUT & 0.089 & 0.424 & 0.681 &  0.001 & 24--96  & 0.435 \\
2024 Park Fire (California)       & USA        & WF & INFL & 0.040 & 0.223 & 0.711 &  0.046 & 24--120 & 0.860 \\
2024 Boise Wildfire (Idaho)       & USA        & WF & INFL & 0.144 & 0.099 & 0.876 &  0.169 & 24--144 & 0.365 \\
\bottomrule
\end{tabular}
\end{table}

\noindent
Abbreviations.  EQ = earthquake; TC = tropical cyclone
(hurricane, typhoon, or tropical storm); FL = flood; WF = wildfire.
EVAC = net near-field population loss ($\dnear < 0$);
INFL = net influx ($\dnear > 0$);
NEUT = near-zero near-field signal ($|\dnear| < 0.01$).


% ============================================================
%  S2  Decay-rate fitting rule
% ============================================================
\section{Decay-rate fitting rule}
\label{sec:si_robustness}

The post-peak power-law exponent $\aemp$ is estimated by OLS on
$\log \Dnorm$ versus $\log t'$ over a window that starts at
$t' = 24$~h to avoid the first post-peak time step (where displacement
is still being resolved at the radial bin level) and terminates at the
first clear demographic rebound, defined as the time index~$i$ at which
$D_{\mathrm{norm},i+1} > 1.05\,D_{\mathrm{norm},i}$.  Fit windows for
each event are listed in Table~\ref{tab:events}.  An alternative
exponential fit produces the same order-of-magnitude spread in recovery
rates (Supplementary Fig.~\ref{fig:S5}), so the qualitative conclusion
that recovery speed varies widely across events does not depend on the
power-law functional form.


% ============================================================
%  S3  PDE Model Details
% ============================================================
\section{PDE model: derivation, Bessel decomposition, and parameter
search}
\label{sec:si_pde}

\subsection*{S3.1 Governing equation}

We model the post-peak evolution of the radial displacement field
$\delta(r,t)$ by a diffusion--relaxation equation in cylindrical
coordinates under radial symmetry:
\begin{equation}
  \frac{\partial\delta}{\partial t}
    = \frac{\Ds}{r}\frac{\partial}{\partial r}
      \!\left(r\,\frac{\partial\delta}{\partial r}\right)
    - k\,\delta, \quad 0 \leq r \leq R,
  \label{si:eq:pde}
\end{equation}
where $\Ds$~[km$^2$\,h$^{-1}$] is a lateral diffusion coefficient
and $k$~[h$^{-1}$] is a uniform first-order relaxation rate.  The
domain radius $R = 200$~km matches the analysis window.

\noindent\textbf{Boundary conditions.}
Neumann (zero-flux) conditions are imposed at both boundaries:
\[
  \left.\frac{\partial\delta}{\partial r}\right|_{r=0} = 0, \qquad
  \left.\frac{\partial\delta}{\partial r}\right|_{r=R} = 0 .
\]

\subsection*{S3.2 Analytical Bessel-series solution}

Because $\delta(r,t)$ satisfies a separable linear PDE with Neumann
boundaries on a disk, the general solution is a superposition of
Bessel-$J_0$ modes~\citep{Crank1975}:
\begin{equation}
  \delta(r,t) = \sum_{n=0}^{N-1} c_n\,J_0\!\left(\frac{\mu_n r}{R}\right)
    e^{-\lambda_n (t-t_{\mathrm{peak}})},
  \label{si:eq:bessel}
\end{equation}
where:
\begin{itemize}
  \item $\mu_n$ are the non-negative zeros of $J_0'(x) = -J_1(x)$
        ($\mu_0 = 0$, $\mu_1 = 3.832$, $\mu_2 = 7.016$, \ldots),
  \item $N = 10$ modes (convergence verified; adding modes
        $n > 10$ changes $\apred$ by $< 0.3\%$),
  \item $\lambda_n = k + \Ds\,(\mu_n/R)^2$ is the modal decay rate.
\end{itemize}

The key physics: higher-order modes (larger $\mu_n$, finer spatial
structure) decay exponentially faster than the fundamental mode.
An initial profile rich in high-order components therefore relaxes
faster in aggregate than a smooth, low-mode profile.

\subsection*{S3.3 Projection of observed profiles onto Bessel basis}

At $t = t_{\mathrm{peak}}$, the expansion coefficients $\{c_n\}$
are determined by $r$-weighted least squares from the observed
10-km-binned radial profile $\hat\delta(r_i)$:
\begin{equation}
  \{c_n\} = \arg\min_{\{c_n\}} \sum_i r_i\,
  \left[\hat\delta(r_i)
        - \sum_{n=0}^{N-1} c_n\,J_0\!\left(\frac{\mu_n r_i}{R}\right)
  \right]^2 .
  \label{si:eq:lsq}
\end{equation}
The weight $r_i$ reflects the area of each annular bin and ensures
that outer bins (larger $r$, more tiles) do not numerically dominate
the fit.

\subsection*{S3.4 Computation of $\apred$}

Rather than tracking $\delta(r,t)$ directly, we use an energy-type
metric:
\begin{equation}
  E(t) = \frac{\int_0^R [\delta(r,t)]^2\,r\,\mathrm{d}r}
               {\int_0^R r\,\mathrm{d}r}
       \;\approx\; \sum_{n=0}^{N-1} c_n^2\,e^{-2\lambda_n
               (t-t_{\mathrm{peak}})}
  \quad\text{(Parseval approximation)}.
  \label{si:eq:energy}
\end{equation}
$E(t)$ is the $L^2$ energy norm, whereas the empirical $D(t)$ is an
$L^1$ mean-absolute-deviation.  Both summarize aggregate relaxation,
but $E(t)$ has an exact Parseval representation that facilitates
interpretation.  We fit $\log E$ vs.\ $\log(t - t_{\mathrm{peak}})$
by OLS over $[1,\,120]$~h to obtain~$\apred$.  The starting point
$1$~h (versus $24$~h for~$\aemp$) is chosen to capture the rapid
initial decay of high-order modes; the model's role is to predict the
ranking of recovery rates, not their absolute magnitude.

\subsection*{S3.5 Global parameter estimation}

Parameters $k$ and $\Ds$ are shared across all 18~events.  We search
a log-uniform $30 \times 30$ grid ($k \in [10^{-4},\,10^{0}]$~h$^{-1}$,
$\Ds \in [10^{-2},\,10^{1}]$~km$^2$\,h$^{-1}$), then refine with a
$30 \times 30$ sub-grid around the coarse optimum (grid spacing reduced
by $10\times$).  The joint optimisation criterion minimises the sum of
three normalised rank statistics:
\begin{equation}
  \mathcal{L}(k, \Ds)
    = \mathrm{rank}\!\left[-\rho_{\mathrm{S}}(\apred,\aemp)\right]
    + \mathrm{rank}\!\left[-r(\apred,\aemp)\right]
    + \mathrm{rank}\!\left[\mathrm{MAE}(\apred,\aemp)\right] ,
  \label{si:eq:criterion}
\end{equation}
where $\rho_{\mathrm{S}}$ denotes Spearman rank correlation,
$r$~Pearson correlation, and MAE is mean absolute error.  The optimal parameters under this criterion are:
\[
  k^* = 0.00418\;\mathrm{h}^{-1}, \quad
  \Ds^* = 0.304\;\mathrm{km}^2\,\mathrm{h}^{-1} .
\]

\subsection*{S3.6 Counterfactual experiments}

Three counterfactual conditions test the necessity of diffusion and
spatial profile structure (Table~\ref{tab:counterfactual}):

\begin{enumerate}
  \item No diffusion ($\Ds = 0$): all modes collapse to the
        same rate~$k$; $E(t) \propto e^{-2kt}$ is profile-independent,
        so $\apred$ is identical for all events and
        $\rho(\apred,\dnear) = 0.0$.
  \item $c_0$-only: only the uniform (zero-wavenumber) mode
        is retained; the spatially heterogeneous structure of the
        profile is erased, again yielding $\rho(\apred,\dnear) = 0.0$.
  \item Shuffled profiles: the observed profile of event~$j$
        is assigned to event~$k$ at random (500~permutations).
        The median $\rho = -0.008$ with 95\% CI $[-0.518,\,+0.519]$
        confirms that the signal requires each event's own
        spatial profile.
\end{enumerate}

\begin{table}[h]
\centering
\caption{Counterfactual experiment results.}
\label{tab:counterfactual}
\small
\begin{tabular}{@{}lccc@{}}
\toprule
Scenario & $\rho(\apred,\dnear)$ & $p$ & 95\% CI \\
\midrule
Baseline PDE (observed) & $-0.356$ & 0.147 & --- \\
No diffusion ($\Ds = 0$) & $0.000$ & 1.000 & --- \\
$c_0$-only & $0.000$ & 1.000 & --- \\
Shuffled profiles & $-0.014$ & --- & $[-0.478,\;+0.445]$ \\
\bottomrule
\end{tabular}
\end{table}

\subsection*{S3.7 Bootstrap robustness of $\rho(\apred, \dnear)$}

To assess sensitivity to the observed radial profiles, we perturb
each event's profile by multiplicative Gaussian noise ($\pm 10\%$)
and refit $\{c_n\}$ for 500~realisations.  The resulting bootstrap
distribution of $\rho(\apred,\dnear)$ remains negatively centred but
is substantially broader than the empirical event-level correlation:
\begin{itemize}
  \item median $\rho = -0.303$,
  \item 95\% CI $= [-0.719,\,+0.270]$,
  \item 87.0\% of realisations yield $\rho < 0$.
\end{itemize}
The bootstrap therefore preserves a negative central tendency, but it
also shows that the mechanistic correlation is sensitive to profile
measurement uncertainty.


% ============================================================
%  S4  Scope: socioeconomic factors
% ============================================================
\section{Scope: socioeconomic factors}
\label{sec:si_socio}

The present study tests a geometric prediction of the diffusion--relaxation
framework and does not jointly model socioeconomic determinants of
recovery (income, institutional quality, infrastructure stock).
Isolating the contribution of these factors would require a harmonised
country- and event-level covariate panel aligned with the current
18-event catalogue; this lies outside the scope of the geometric
analysis reported here.


% ============================================================
%  S5  Distance-Segment Analysis
% ============================================================
\section{Distance-segment robustness of the shape--rate correlation}
\label{sec:si_distance}

To test whether the shape--rate signal is driven by a specific radial
zone, we compute $\rho(\aemp,\bar\delta_{\mathrm{zone}})$ separately
for near-field ($r < 50$~km), mid-field (50--100~km), and far-field
($100$--$200$~km) segments.  $\bar\delta_{\mathrm{zone}}$ is the
spatial mean of the displacement field within each zone at
$t_{\mathrm{peak}}$.

\begin{itemize}
  \item Near-field ($r < 50$~km, $n = 18$):
        $\rho = -0.672$, $p = 0.002$.
  \item Mid-field (50--100~km, $n = 18$):
        $\rho = -0.686$, $p = 0.002$.  This segment remains strongly
        associated with recovery speed, consistent with a coherent
        mesoscale evacuation shadow around fast-recovering events.
  \item Far-field ($100$--$200$~km, $n = 18$):
        $\rho = -0.556$, $p = 0.017$.  The association weakens with
        distance but remains negative, indicating that the clearest
        shape--rate signal originates near the event centre.
\end{itemize}

The primary predictor $\dnear$ (near-field average) is chosen because
it maximises sample coverage while remaining physically interpretable
as the net population flux near the disaster epicenter.


% ============================================================
%  S6  Additional Robustness Results
% ============================================================
\section{Additional robustness results}
\label{sec:si_additional}

\subsection*{S6.1 Leave-one-out jackknife}

Jackknife resampling (leave-one-out, $n=18$ iterations) yields a
95\% CI for $\rho(\aemp,\dnear)$ that does not cross zero: the
leave-one-out replicate distribution spans $[-0.743,\,-0.650]$,
confirming that no single event drives the correlation.

\subsection*{S6.2 Residual displacement vs.\ decay rate}

The residual displacement $D_\infty$ (mean of normalised $D(t)$
over the post-peak tail) is strongly anticorrelated with~$\aemp$:
$\rho(\aemp, D_\infty) = -0.674$, $p = 0.002$.  Events with
faster decay also reach lower residual levels, consistent with the
interpretation that $\aemp$ captures genuine recovery rather than a
transient fluctuation.

\subsection*{S6.3 Excluded candidate predictors}

Two additional candidate predictors of $\aemp$ were tested but produced
$|\rho| < 0.15$ and $p > 0.5$ in bivariate analysis: the recovery
gradient $\mathrm{d}D/\mathrm{d}t$ at $t_{\mathrm{peak}}$, and the
observation window length (time between $t_{\mathrm{peak}}$ and the
last retained post-peak step). Neither enters the main analysis.

\subsection*{S6.4 Hazard-type grouping test}

To assess whether recovery speed varies systematically by hazard type,
we group the 18-event primary event-level sample into four broad
categories: earthquake ($n = 1$), tropical cyclone ($n = 7$; hurricanes,
typhoons, and tropical storms combined), flood ($n = 7$), and wildfire
($n = 3$). A Kruskal--Wallis rank-sum test across these four groups
gives $H = 3.80$, $p = 0.28$, indicating no significant difference in
recovery speed across broad hazard categories. As a sensitivity check,
excluding the singleton earthquake group and retaining only categories
with $n \geq 2$ (tropical cyclone, flood, wildfire) gives
$H = 3.44$, $p = 0.18$. In both cases, within-group variation in
$\aemp$ exceeds between-group variation, indicating that hazard type
alone does not explain the observed diversity of recovery speeds.


% ============================================================
%  Supplementary Figures
% ============================================================
\section{Supplementary figures}
\label{sec:si_figures}

\begin{figure}[h]
\centering
\includegraphics[width=0.95\linewidth]{%
  figures_supp/figS1_individual_panels}
\caption{Individual $D(t)$ curves for all selected events.  Each panel
shows aggregated displacement and event-level fit diagnostics
($\alpha$, $R^2$).}
\label{fig:S1}
\end{figure}

\begin{figure}[h]
\centering
\includegraphics[width=0.75\linewidth]{%
  figures_supp/figS2_jackknife}
\caption{Leave-one-out jackknife distribution of
$\rho(\aemp,\dnear)$, with full-sample baseline marked.}
\label{fig:S2}
\end{figure}

\begin{figure}[h]
\centering
\includegraphics[width=0.60\linewidth]{%
  figures_supp/figS3_alpha_vs_dinf}
\caption{$\alpha$ versus residual displacement $D_\infty$.
This panel is retained as a supplementary validation of decay-rate
interpretability.}
\label{fig:S3}
\end{figure}

\begin{figure}[h]
\centering
\includegraphics[width=0.95\linewidth]{%
  figures_supp/figS4_profile_gallery}
\caption{Peak-time radial displacement profiles
$\delta(r,t_{\mathrm{peak}})$ for selected events.}
\label{fig:S4}
\end{figure}

\begin{figure}[h]
\centering
\includegraphics[width=0.70\linewidth]{%
  figures_supp/figS5_gao_delta_bic}
\caption{Model-selection support figure:
$\Delta$BIC (power law minus exponential) by event.}
\label{fig:S5}
\end{figure}

\bibliographystyle{pnas2009}
\bibliography{references}

\end{document}
